\section{Working conventions before the first calculation}
\label{app:conventions}
% BEGIN CURATED SUBJECT INDEX
\index{GNS construction}
\index{boundary condition!incoming and outgoing}
\index{antiresonance}
\index{quasinormal mode}
\index{horizon boundary condition}
% END CURATED SUBJECT INDEX

We use time dependence $\rme^{-\rmi Et/\hbar}$ in quantum mechanics and
$\rme^{-\rmi\omega t}$ in black-hole perturbation theory.  A decaying pole has
\begin{equation}
  \im \resonantE<0,
  \qquad
  \im\vsub{\omega}{QNM}<0 .
  \label{eq:app-decay-signs}
\end{equation}
For $E=\hbar^2k^2/(2m)$, the physical scattering rim has $k>0$ for
$E>0$.  Continuing through the positive-$E$ cut changes the sign of the square
root.  A resonance is conventionally represented by $\re k>0$ and
$\im k<0$; its time-reversed antiresonance is obtained using the reality
symmetry of the potential.

The one-dimensional probability current is
\begin{equation}
  j[\psi]
  =\frac{\hbar}{2m\rmi}
  \left(\psi^*\partial_x\psi
  -\psi\partial_x\psi^*\right) .
  \label{eq:probability-current}
\end{equation}
For real $k>0$, $\rme^{+\rmi kx}$ carries positive current and
$\rme^{-\rmi kx}$ carries negative current.  At a complex resonance momentum,
the current of a non-normalizable stationary wave is not a conserved
probability flux in the ordinary Hilbert sense; outgoing/incoming labels are
defined by analytic continuation from real $k$.

In black-hole coordinates, the group direction is most transparently read
from null combinations.  Near a horizon, a factor
$\rme^{-\rmi\omega(t+r_*)}$ is ingoing, while
$\rme^{-\rmi\omega(t-r_*)}$ is outgoing.  This gives the signs in
Eq.~\eqref{eq:qnm-boundary-flat}.
